% --- Archon physics formalization source begin ---
% archon:physics
% archon:covers PhyXMiniProblems/problem_phyx_mini_0609.lean
% archon:source-report reports/phyx_mini/problem_phyx_mini_0609.source.json

\chapter{Physics problem phyx\_mini\_0609}
\label{ch:PhyXMiniProblems_problem_phyx_mini_0609}

\paragraph{Problem source.}
\#\# Physical scenario

In a laboratory, a rectangular loop of wire surrounds the origin in the \$xz\$-plane, with extent \$H\$ in the \$z\$-direction and extent \$L\$ in the \$x\$-direction (figure). The loop carries current \$I\$ in the counterclockwise direction as viewed from the positive \$y\$-axis. A magnetic field \$\textbackslash{}vec\{B\} = B \textbackslash{}hat\{k\}\$ is present. The electromagnetic field is \$\textbackslash{}vec\{E\} \textbackslash{}perp \textbackslash{}vec\{B\} = (0, \textbackslash{}vec\{B\})\$ in an inertial frame \$S\$, and then in another frame \$S'\$ moving at velocity \$\textbackslash{}vec\{v\}\$ relative to \$S\$.

\#\# Question

What is the component of the magnetic field \$\textbackslash{}vec\{B\}'\$, perpendicular to \$\textbackslash{}vec\{v\}\$?

\#\# Answer choices

- A: A: \$\textbackslash{}mathbf\{B\}\_\textbackslash{}gamma\$

- B: B: \$\textbackslash{}mathbf\{B\}\_\textbackslash{}alpha\$

- C: C: \$\textbackslash{}mathbf\{B\}\_\textbackslash{}beta\$

- D: D: \$\textbackslash{}mathbf\{B\}\_\textbackslash{}delta\$

\#\# Figure caption (auxiliary; use the image as primary evidence)

The image depicts a three-dimensional illustration of a rectangular loop in a magnetic field, with several vectors and labels:

1. **Rectangular Loop**: Represented in a light brown color, placed in a 3D coordinate system (x, y, z).

2. **Coordinate Axes**: The x-axis is horizontal, the y-axis is perpendicular to it, and the z-axis is vertical.

3. **Vectors**:
   - \textbackslash{}(\textbackslash{}vec\{B\}\textbackslash{}) (blue arrow pointing upwards along the z-axis): Represents the magnetic field.
   - \textbackslash{}(\textbackslash{}vec\{\textbackslash{}mu\}\textbackslash{}) (red arrow pointing to the right along the y-axis): Indicates the magnetic moment.
   - \textbackslash{}(\textbackslash{}vec\{I\}\textbackslash{}) (purple arrow along the right side of the loop): Represents the current direction through the loop.
   - \textbackslash{}(\textbackslash{}vec\{\textbackslash{}tau\}\textbackslash{}) (curved arrow along the x-axis, denoting a counter-clockwise direction): Signifying torque.

4. **Dimensions**:
   - \textbackslash{}(L\textbackslash{}): Length of the side of the loop along the z-axis.
   - \textbackslash{}(H\textbackslash{}): Height of the side of the loop along the x-axis.

The image illustrates the interaction of a current-carrying loop with a magnetic field, indicating torque, field direction, and current flow.

\#\# Dataset metadata

Relativity; Physical Model Grounding Reasoning, Spatial Relation Reasoning

\paragraph{Recorded answer/context.}
A: A: \$\textbackslash{}mathbf\{B\}\_\textbackslash{}gamma\$

\paragraph{Figure/image path.}
/root/proposal\_for\_physic/science-mango/phyx\_mini\_run/phyx\_data/test\_image/609.png

\paragraph{Formalization target.}
create a compiling Lean file with sorry bodies at `PhyXMiniProblems/problem\_phyx\_mini\_0609.lean`. The Lean declarations must preserve the physical quantities, dimensions or dimensional roles, figure labels, governing-law hypotheses, and final relation expressed by this problem.
Use Mathlib/Physlib names found through LeanExplore where available. If a physics API is missing, introduce faithful local abstractions rather than scalar placeholder aliases.

\begin{theorem}[Physics formalization target]
\label{thm:physics:phyx_mini_0609:target}
\lean{PhyXMiniProblems.ProblemPhyXMini0609.perpendicularMagneticFieldInMovingFrame_eq_gamma_smul}
\uses{def:physics:phyx-mini-0609:phyxminiproblems-problemphyxmini0609-perpendicularcomponent, def:physics:phyx-mini-0609:phyxminiproblems-problemphyxmini0609-movingframemagneticfieldsetup, def:physics:phyx-mini-0609:phyxminiproblems-problemphyxmini0609-magneticvectorat, def:physics:phyx-mini-0609:phyxminiproblems-problemphyxmini0609-matchesmovingframemagneticfieldscenario, def:physics:phyx-mini-0609:phyxminiproblems-problemphyxmini0609-hasphysicalmovingframeparameters, def:physics:phyx-mini-0609:phyxminiproblems-problemphyxmini0609-matchessuppliedrectangularcurrentloopfigure, def:physics:phyx-mini-0609:phyxminiproblems-problemphyxmini0609-satisfieslorentzelectromagneticfieldtransformation}
The assigned autoformalize agent should translate this physics problem into Lean declarations in the covered file, with theorem and lemma proof bodies written as `by sorry`.
\end{theorem}
\begin{proof}
This is an autoformalization task, not a proof task. Produce statements that can later be proved by the physics prover without weakening the physical model.
\end{proof}

% --- Archon declaration topology begin ---
\section{Declaration topology}
\begin{definition}[Spatial Vector]
  \label{def:physics:phyx-mini-0609:phyxminiproblems-problemphyxmini0609-spatialvector}
  \lean{PhyXMiniProblems.ProblemPhyXMini0609.SpatialVector}
  A three-dimensional spatial vector in Physlib's Euclidean space
\end{definition}
\begin{proof}
  This is the declared model definition of Spatial Vector.
\end{proof}
\begin{definition}[electric Current Dimension]
  \label{def:physics:phyx-mini-0609:phyxminiproblems-problemphyxmini0609-electriccurrentdimension}
  \lean{PhyXMiniProblems.ProblemPhyXMini0609.electricCurrentDimension}
  Electric current has physical dimension charge per unit time
\end{definition}
\begin{proof}
  This is the declared model definition of electric Current Dimension.
\end{proof}
\begin{definition}[magnetic Moment Dimension]
  \label{def:physics:phyx-mini-0609:phyxminiproblems-problemphyxmini0609-magneticmomentdimension}
  \lean{PhyXMiniProblems.ProblemPhyXMini0609.magneticMomentDimension}
  \uses{def:physics:phyx-mini-0609:phyxminiproblems-problemphyxmini0609-electriccurrentdimension}
  A magnetic dipole moment has dimension current times area
\end{definition}
\begin{proof}
  This is the declared model definition of magnetic Moment Dimension.
\end{proof}
\begin{definition}[torque Dimension]
  \label{def:physics:phyx-mini-0609:phyxminiproblems-problemphyxmini0609-torquedimension}
  \lean{PhyXMiniProblems.ProblemPhyXMini0609.torqueDimension}
  Torque has the same physical dimension as energy
\end{definition}
\begin{proof}
  This is the declared model definition of torque Dimension.
\end{proof}
\begin{definition}[Length Quantity]
  \label{def:physics:phyx-mini-0609:phyxminiproblems-problemphyxmini0609-lengthquantity}
  \lean{PhyXMiniProblems.ProblemPhyXMini0609.LengthQuantity}
  A nonnegative, unit-independent physical length
\end{definition}
\begin{proof}
  This is the declared model definition of Length Quantity.
\end{proof}
\begin{definition}[Current Quantity]
  \label{def:physics:phyx-mini-0609:phyxminiproblems-problemphyxmini0609-currentquantity}
  \lean{PhyXMiniProblems.ProblemPhyXMini0609.CurrentQuantity}
  \uses{def:physics:phyx-mini-0609:phyxminiproblems-problemphyxmini0609-electriccurrentdimension}
  A nonnegative, unit-independent electric-current magnitude
\end{definition}
\begin{proof}
  This is the declared model definition of Current Quantity.
\end{proof}
\begin{definition}[Magnetic Moment Quantity]
  \label{def:physics:phyx-mini-0609:phyxminiproblems-problemphyxmini0609-magneticmomentquantity}
  \lean{PhyXMiniProblems.ProblemPhyXMini0609.MagneticMomentQuantity}
  \uses{def:physics:phyx-mini-0609:phyxminiproblems-problemphyxmini0609-spatialvector, def:physics:phyx-mini-0609:phyxminiproblems-problemphyxmini0609-magneticmomentdimension}
  A unit-independent magnetic-moment vector
\end{definition}
\begin{proof}
  This is the declared model definition of Magnetic Moment Quantity.
\end{proof}
\begin{definition}[Torque Quantity]
  \label{def:physics:phyx-mini-0609:phyxminiproblems-problemphyxmini0609-torquequantity}
  \lean{PhyXMiniProblems.ProblemPhyXMini0609.TorqueQuantity}
  \uses{def:physics:phyx-mini-0609:phyxminiproblems-problemphyxmini0609-spatialvector, def:physics:phyx-mini-0609:phyxminiproblems-problemphyxmini0609-torquedimension}
  A unit-independent torque vector
\end{definition}
\begin{proof}
  This is the declared model definition of Torque Quantity.
\end{proof}
\begin{definition}[nonnegative SIReadout]
  \label{def:physics:phyx-mini-0609:phyxminiproblems-problemphyxmini0609-nonnegativesireadout}
  \lean{PhyXMiniProblems.ProblemPhyXMini0609.nonnegativeSIReadout}
  Coherent-SI scalar readout of a nonnegative dimensionful quantity
\end{definition}
\begin{proof}
  This is the declared model definition of nonnegative SIReadout.
\end{proof}
\begin{definition}[vector SIReadout]
  \label{def:physics:phyx-mini-0609:phyxminiproblems-problemphyxmini0609-vectorsireadout}
  \lean{PhyXMiniProblems.ProblemPhyXMini0609.vectorSIReadout}
  \uses{def:physics:phyx-mini-0609:phyxminiproblems-problemphyxmini0609-spatialvector}
  Coherent-SI vector readout of a dimensionful spatial vector
\end{definition}
\begin{proof}
  This is the declared model definition of vector SIReadout.
\end{proof}
\begin{definition}[Coordinate Axis]
  \label{def:physics:phyx-mini-0609:phyxminiproblems-problemphyxmini0609-coordinateaxis}
  \lean{PhyXMiniProblems.ProblemPhyXMini0609.CoordinateAxis}
  The three coordinate axes appearing in the prose and the figure
\end{definition}
\begin{proof}
  This is the declared model definition of Coordinate Axis.
\end{proof}
\begin{definition}[axis Vector]
  \label{def:physics:phyx-mini-0609:phyxminiproblems-problemphyxmini0609-axisvector}
  \lean{PhyXMiniProblems.ProblemPhyXMini0609.axisVector}
  \uses{def:physics:phyx-mini-0609:phyxminiproblems-problemphyxmini0609-spatialvector, def:physics:phyx-mini-0609:phyxminiproblems-problemphyxmini0609-coordinateaxis}
  Unit coordinate vector associated with a labelled axis
\end{definition}
\begin{proof}
  This is the declared model definition of axis Vector.
\end{proof}
\begin{definition}[Points Along Positive Axis]
  \label{def:physics:phyx-mini-0609:phyxminiproblems-problemphyxmini0609-pointsalongpositiveaxis}
  \lean{PhyXMiniProblems.ProblemPhyXMini0609.PointsAlongPositiveAxis}
  \uses{def:physics:phyx-mini-0609:phyxminiproblems-problemphyxmini0609-spatialvector, def:physics:phyx-mini-0609:phyxminiproblems-problemphyxmini0609-coordinateaxis, def:physics:phyx-mini-0609:phyxminiproblems-problemphyxmini0609-axisvector}
  A vector points strictly along the positive direction of a labelled axis
\end{definition}
\begin{proof}
  This is the declared model definition of Points Along Positive Axis.
\end{proof}
\begin{definition}[spatial Cross]
  \label{def:physics:phyx-mini-0609:phyxminiproblems-problemphyxmini0609-spatialcross}
  \lean{PhyXMiniProblems.ProblemPhyXMini0609.spatialCross}
  \uses{def:physics:phyx-mini-0609:phyxminiproblems-problemphyxmini0609-spatialvector}
  Mathlib's ordinary three-dimensional cross product, transported to SpatialVector
\end{definition}
\begin{proof}
  This is the declared model definition of spatial Cross.
\end{proof}
\begin{definition}[perpendicular Component]
  \label{def:physics:phyx-mini-0609:phyxminiproblems-problemphyxmini0609-perpendicularcomponent}
  \lean{PhyXMiniProblems.ProblemPhyXMini0609.perpendicularComponent}
  \uses{def:physics:phyx-mini-0609:phyxminiproblems-problemphyxmini0609-spatialvector}
  The component perpendicular to a unit direction axis. Unit normalization is an explicit admissibility premise below rather than being hidden in this geometric definition
\end{definition}
\begin{proof}
  This is the declared model definition of perpendicular Component.
\end{proof}
\begin{definition}[Inertial Frame]
  \label{def:physics:phyx-mini-0609:phyxminiproblems-problemphyxmini0609-inertialframe}
  \lean{PhyXMiniProblems.ProblemPhyXMini0609.InertialFrame}
  The laboratory frame and the moving inertial frame
\end{definition}
\begin{proof}
  This is the declared model definition of Inertial Frame.
\end{proof}
\begin{definition}[Coordinate Plane]
  \label{def:physics:phyx-mini-0609:phyxminiproblems-problemphyxmini0609-coordinateplane}
  \lean{PhyXMiniProblems.ProblemPhyXMini0609.CoordinatePlane}
  Coordinate planes relevant to the rectangular loop
\end{definition}
\begin{proof}
  This is the declared model definition of Coordinate Plane.
\end{proof}
\begin{definition}[Circulation Sense]
  \label{def:physics:phyx-mini-0609:phyxminiproblems-problemphyxmini0609-circulationsense}
  \lean{PhyXMiniProblems.ProblemPhyXMini0609.CirculationSense}
  Sense in which the current circulates when viewed from a positive axis
\end{definition}
\begin{proof}
  This is the declared model definition of Circulation Sense.
\end{proof}
\begin{definition}[Dimension Label]
  \label{def:physics:phyx-mini-0609:phyxminiproblems-problemphyxmini0609-dimensionlabel}
  \lean{PhyXMiniProblems.ProblemPhyXMini0609.DimensionLabel}
  Literal dimension labels printed in the supplied image
\end{definition}
\begin{proof}
  This is the declared model definition of Dimension Label.
\end{proof}
\begin{definition}[Figure Arrow Label]
  \label{def:physics:phyx-mini-0609:phyxminiproblems-problemphyxmini0609-figurearrowlabel}
  \lean{PhyXMiniProblems.ProblemPhyXMini0609.FigureArrowLabel}
  Vector and current labels printed beside arrows in the supplied image
\end{definition}
\begin{proof}
  This is the declared model definition of Figure Arrow Label.
\end{proof}
\begin{definition}[Positive Axis Direction]
  \label{def:physics:phyx-mini-0609:phyxminiproblems-problemphyxmini0609-positiveaxisdirection}
  \lean{PhyXMiniProblems.ProblemPhyXMini0609.PositiveAxisDirection}
  Positive coordinate direction assigned to a straight vector arrow or axis
\end{definition}
\begin{proof}
  This is the declared model definition of Positive Axis Direction.
\end{proof}
\begin{definition}[Rectangular Current Loop Figure]
  \label{def:physics:phyx-mini-0609:phyxminiproblems-problemphyxmini0609-rectangularcurrentloopfigure}
  \lean{PhyXMiniProblems.ProblemPhyXMini0609.RectangularCurrentLoopFigure}
  \uses{def:physics:phyx-mini-0609:phyxminiproblems-problemphyxmini0609-coordinateaxis, def:physics:phyx-mini-0609:phyxminiproblems-problemphyxmini0609-coordinateplane, def:physics:phyx-mini-0609:phyxminiproblems-problemphyxmini0609-circulationsense, def:physics:phyx-mini-0609:phyxminiproblems-problemphyxmini0609-dimensionlabel, def:physics:phyx-mini-0609:phyxminiproblems-problemphyxmini0609-figurearrowlabel, def:physics:phyx-mini-0609:phyxminiproblems-problemphyxmini0609-positiveaxisdirection}
  Literal and qualitative evidence visible in image 609
\end{definition}
\begin{proof}
  This is the declared model definition of Rectangular Current Loop Figure.
\end{proof}
\begin{definition}[Moving Frame Magnetic Field Setup]
  \label{def:physics:phyx-mini-0609:phyxminiproblems-problemphyxmini0609-movingframemagneticfieldsetup}
  \lean{PhyXMiniProblems.ProblemPhyXMini0609.MovingFrameMagneticFieldSetup}
  \uses{def:physics:phyx-mini-0609:phyxminiproblems-problemphyxmini0609-spatialvector, def:physics:phyx-mini-0609:phyxminiproblems-problemphyxmini0609-lengthquantity, def:physics:phyx-mini-0609:phyxminiproblems-problemphyxmini0609-currentquantity, def:physics:phyx-mini-0609:phyxminiproblems-problemphyxmini0609-magneticmomentquantity, def:physics:phyx-mini-0609:phyxminiproblems-problemphyxmini0609-torquequantity, def:physics:phyx-mini-0609:phyxminiproblems-problemphyxmini0609-inertialframe, def:physics:phyx-mini-0609:phyxminiproblems-problemphyxmini0609-coordinateplane, def:physics:phyx-mini-0609:phyxminiproblems-problemphyxmini0609-circulationsense, def:physics:phyx-mini-0609:phyxminiproblems-problemphyxmini0609-rectangularcurrentloopfigure}
  The loop quantities, frame-dependent electromagnetic fields, two coordinate descriptions of one observation event, and relative-motion parameters are independent data. The magnetic field in SPrime is deliberately not defined from γ or from an answer choice
\end{definition}
\begin{proof}
  This is the declared model definition of Moving Frame Magnetic Field Setup.
\end{proof}
\begin{definition}[electric Vector At]
  \label{def:physics:phyx-mini-0609:phyxminiproblems-problemphyxmini0609-electricvectorat}
  \lean{PhyXMiniProblems.ProblemPhyXMini0609.electricVectorAt}
  \uses{def:physics:phyx-mini-0609:phyxminiproblems-problemphyxmini0609-spatialvector, def:physics:phyx-mini-0609:phyxminiproblems-problemphyxmini0609-inertialframe, def:physics:phyx-mini-0609:phyxminiproblems-problemphyxmini0609-movingframemagneticfieldsetup}
  Electric-field vector at the selected event in a specified frame
\end{definition}
\begin{proof}
  This is the declared model definition of electric Vector At.
\end{proof}
\begin{definition}[magnetic Vector At]
  \label{def:physics:phyx-mini-0609:phyxminiproblems-problemphyxmini0609-magneticvectorat}
  \lean{PhyXMiniProblems.ProblemPhyXMini0609.magneticVectorAt}
  \uses{def:physics:phyx-mini-0609:phyxminiproblems-problemphyxmini0609-spatialvector, def:physics:phyx-mini-0609:phyxminiproblems-problemphyxmini0609-inertialframe, def:physics:phyx-mini-0609:phyxminiproblems-problemphyxmini0609-movingframemagneticfieldsetup}
  Magnetic-field vector at the selected event in a specified frame
\end{definition}
\begin{proof}
  This is the declared model definition of magnetic Vector At.
\end{proof}
\begin{definition}[Matches Moving Frame Magnetic Field Scenario]
  \label{def:physics:phyx-mini-0609:phyxminiproblems-problemphyxmini0609-matchesmovingframemagneticfieldscenario}
  \lean{PhyXMiniProblems.ProblemPhyXMini0609.MatchesMovingFrameMagneticFieldScenario}
  \uses{def:physics:phyx-mini-0609:phyxminiproblems-problemphyxmini0609-vectorsireadout, def:physics:phyx-mini-0609:phyxminiproblems-problemphyxmini0609-axisvector, def:physics:phyx-mini-0609:phyxminiproblems-problemphyxmini0609-pointsalongpositiveaxis, def:physics:phyx-mini-0609:phyxminiproblems-problemphyxmini0609-movingframemagneticfieldsetup, def:physics:phyx-mini-0609:phyxminiproblems-problemphyxmini0609-electricvectorat, def:physics:phyx-mini-0609:phyxminiproblems-problemphyxmini0609-magneticvectorat}
  Prose-level assignments, including (E, B) = (0, B zHat) in frame S
\end{definition}
\begin{proof}
  This is the declared model definition of Matches Moving Frame Magnetic Field Scenario.
\end{proof}
\begin{definition}[Has Physical Moving Frame Parameters]
  \label{def:physics:phyx-mini-0609:phyxminiproblems-problemphyxmini0609-hasphysicalmovingframeparameters}
  \lean{PhyXMiniProblems.ProblemPhyXMini0609.HasPhysicalMovingFrameParameters}
  \uses{def:physics:phyx-mini-0609:phyxminiproblems-problemphyxmini0609-nonnegativesireadout, def:physics:phyx-mini-0609:phyxminiproblems-problemphyxmini0609-movingframemagneticfieldsetup}
  Positivity, subluminality, and normalization conditions for the setup
\end{definition}
\begin{proof}
  This is the declared model definition of Has Physical Moving Frame Parameters.
\end{proof}
\begin{definition}[Matches Supplied Rectangular Current Loop Figure]
  \label{def:physics:phyx-mini-0609:phyxminiproblems-problemphyxmini0609-matchessuppliedrectangularcurrentloopfigure}
  \lean{PhyXMiniProblems.ProblemPhyXMini0609.MatchesSuppliedRectangularCurrentLoopFigure}
  \uses{def:physics:phyx-mini-0609:phyxminiproblems-problemphyxmini0609-movingframemagneticfieldsetup}
  Primary-image evidence. The raster places H along the vertical z extent and L along the slanted x extent, resolving the reversed assignment in the auxiliary caption in favor of the image and scenario prose
\end{definition}
\begin{proof}
  This is the declared model definition of Matches Supplied Rectangular Current Loop Figure.
\end{proof}
\begin{definition}[Satisfies Lorentz Electromagnetic Field Transformation]
  \label{def:physics:phyx-mini-0609:phyxminiproblems-problemphyxmini0609-satisfieslorentzelectromagneticfieldtransformation}
  \lean{PhyXMiniProblems.ProblemPhyXMini0609.SatisfiesLorentzElectromagneticFieldTransformation}
  \uses{def:physics:phyx-mini-0609:phyxminiproblems-problemphyxmini0609-spatialcross, def:physics:phyx-mini-0609:phyxminiproblems-problemphyxmini0609-perpendicularcomponent, def:physics:phyx-mini-0609:phyxminiproblems-problemphyxmini0609-movingframemagneticfieldsetup, def:physics:phyx-mini-0609:phyxminiproblems-problemphyxmini0609-electricvectorat, def:physics:phyx-mini-0609:phyxminiproblems-problemphyxmini0609-magneticvectorat}
  For the convention that SPrime moves with velocity v relative to S, the general transverse magnetic-field law at corresponding events is B'\_perp = γ(β) (B\_perp - (1/c²) (v × E)). The electric mixing term is retained, so the desired zero-electric-field specialization is not assumed here
\end{definition}
\begin{proof}
  This is the declared model definition of Satisfies Lorentz Electromagnetic Field Transformation.
\end{proof}
\begin{definition}[Answer Choice]
  \label{def:physics:phyx-mini-0609:phyxminiproblems-problemphyxmini0609-answerchoice}
  \lean{PhyXMiniProblems.ProblemPhyXMini0609.AnswerChoice}
  The four answer-choice labels
\end{definition}
\begin{proof}
  This is the declared model definition of Answer Choice.
\end{proof}
\begin{definition}[Field Subscript]
  \label{def:physics:phyx-mini-0609:phyxminiproblems-problemphyxmini0609-fieldsubscript}
  \lean{PhyXMiniProblems.ProblemPhyXMini0609.FieldSubscript}
  Subscript printed on the magnetic-field symbol in an answer choice
\end{definition}
\begin{proof}
  This is the declared model definition of Field Subscript.
\end{proof}
\begin{definition}[displayed Subscript]
  \label{def:physics:phyx-mini-0609:phyxminiproblems-problemphyxmini0609-displayedsubscript}
  \lean{PhyXMiniProblems.ProblemPhyXMini0609.displayedSubscript}
  \uses{def:physics:phyx-mini-0609:phyxminiproblems-problemphyxmini0609-answerchoice, def:physics:phyx-mini-0609:phyxminiproblems-problemphyxmini0609-fieldsubscript}
  Literal mapping from the printed choices to their displayed subscripts
\end{definition}
\begin{proof}
  This is the declared model definition of displayed Subscript.
\end{proof}
% --- Archon declaration topology end ---
% --- Archon physics formalization source end ---
